\chapter{أنواع العصبونات بحسب وظيفتها الكهربائية}
\chaptermark{العصبونات أجهزةً}
\label{sec:eetypes}

صنّفنا العصبونات في الفصل~\ref{sec:neurontypes} بحسب الناقل العصبي الذي يطلقه مخرجها. أما مهندس الإلكترونيات فكان سيصنّفها بطريقة أخرى: بحسب ما يفعله العصبون بإشارته، أي بحسب الجهاز الذي يعمل عمله. الجهاز الرئيسي في الكتاب نعرفه منذ الفصل~\ref{sec:glutcomputer}: العصبون المكامل التسريبي~\eqref{eq:glutneuron}، أي دارة RC مع مقارِن. لكن في الدماغ أجهزة أخرى. جُمعت في الجدول~\ref{tab:eetypes} في أربع مجموعات، وقد مرّ كل منها تقريبًا في الكتاب.

\begin{center}
\footnotesize
\setlength{\tabcolsep}{3pt}
\renewcommand{\arraystretch}{1.15}
\begin{longtable}{>{\raggedright}p{3.1cm}>{\raggedright}p{3.3cm}>{\raggedright}p{4.7cm}>{\raggedright\arraybackslash}p{2.4cm}}
\caption{العصبونات أجهزةً: الاسم، والنظير في الإلكترونيات، وما يفعله الجهاز، وأين يرد في الكتاب.}\label{tab:eetypes}\\
\toprule
العصبون & الجهاز & ماذا يفعل & أين في الكتاب \\
\midrule
\endfirsthead
\toprule
العصبون & الجهاز & ماذا يفعل & أين في الكتاب \\
\midrule
\endhead
\bottomrule
\endfoot
\multicolumn{4}{l}{\emph{المرشحات والمكاملات}} \\
العصبون المكامل التسريبي & مكامل RC مع مقارِن & يجمع المداخل في نافذة $\tau$ ويطلق عند العتبة & الفصل~\ref{sec:glutcomputer}، \eqref{eq:glutneuron} \\
كاشف التزامن & بوابة AND بنافذة زمنية & لا يطلق إلا إذا وصلت المداخل معًا تقريبًا؛ $\tau$ صغيرة جدًا & الفصول~\ref{sec:glutcomputer}، \ref{sec:code}، \eqref{eq:window} \\
العصبون الطوري & مرشح تمرير عالٍ، كاشف حافة & يستجيب لتغير المدخل ثم يصمت & الفصل~\ref{sec:sensors} \\
العصبون المتكيّف & مرشح تمرير عالٍ مع تحكم آلي في الكسب & يهبط التردد مع المدخل الثابت & الفصل~\ref{sec:sensors}، \eqref{eq:nakarushton} \\
الرنّان & دارة رنين، مرشح تمرير نطاقي & أقوى استجابته لمدخل بتردده الخاص & القسم~\ref{sec:resonator} \\
المكامل غير التسريبي & مكامل بمضخم عمليات & يحوّل السرعة إلى موضع ويحفظه & القسم~\ref{sec:perfectint} \\
\midrule
\multicolumn{4}{l}{\emph{المحوّلات}} \\
العصبون التوتري & محوّل جهد إلى تردد & يتبع التردد المدخلَ خطيًا ولا يكاد يتعب & الفصل~\ref{sec:code} \\
العصبون المتدرج (بلا نبضات) & مضخم تماثلي، عازل & ينقل جهدًا متصلًا إلى مسافة قصيرة & الفصل~\ref{sec:sensors} \\
العصبون المقوِّم & مقوّم نصف موجة & التردد لا يكون سالبًا، $[u]_+$ & الفصلان~\ref{sec:glutcomputer}، \ref{sec:gaba} \\
زوج ”تشغيل/إطفاء“ & زوج مقوّمات، شفرة ثنائية السلك & أحدهما يستجيب لـ”أكثر“ والآخر لـ”أقل“ & الفصل~\ref{sec:glutcomputer}، \eqref{eq:dualrail} \\
عصبون بتثبيط تحويلي & مقسِّم تماثلي، ضبط الكسب & يقسم المدخل بدل أن يطرح منه & الفصل~\ref{sec:gaba}، \eqref{eq:shunt} \\
\midrule
\multicolumn{4}{l}{\emph{المولّدات والزمن}} \\
ناظمة الإيقاع & مذبذب استرخائي، هزّاز متعدد & يُصدر بنفسه نبضات بتردد ثابت & الفصلان~\ref{sec:serocomputer}، \ref{sec:dofa} \\
ناظمة إيقاع لها مدخل & مذبذب متحكَّم فيه بالجهد & إيقاع خاص يسرّعه المدخل أو يبطئه & الفصل~\ref{sec:chemoutput} \\
عصبون الرشقات & مولّد رشقات مبوَّب & يُصدر رشقات من نبضات متقاربة، أي ”شَرطات“ & الفصل~\ref{sec:code}، \eqref{eq:burst} \\
عصبون مقيَّد بالطور & كاشف طور، حلقة مقفلة الطور & يطلق عند طور محدد من الموجة الداخلة & الفصلان~\ref{sec:sensors}، \ref{sec:chemoutput} \\
محوار بتأخير & خط تأخير & يؤخر الإشارة زمنًا محددًا & الفصل~\ref{sec:glutcomputer}، الشكل~\ref{fig:itd} \\
\midrule
\multicolumn{4}{l}{\emph{الذاكرة والتبديل}} \\
العصبون ثنائي الاستقرار & قادح شميت، ماسك & مدخل قصير ينقله إلى حالة طويلة الأمد & القسم~\ref{sec:bistable} \\
عصبون بفروع تغصنية & مُضاعِف إرسال، شبكة صغيرة متعددة الطبقات & الفرع لا يمرّر المدخل إلا إذا وُجد مدخل سماح & الفصل~\ref{sec:synthesis} \\
\end{longtable}
\end{center}

تحفّظ واحد أهم من الجدول كله: هذه ليست خلايا مختلفة، بل \emph{أوضاع} مختلفة. العصبون نفسه قد يكون مكاملًا مع المدخل البطيء وكاشف تزامن حين يقصّر التثبيطُ $\tau$ الخاصة به (الفصل~\ref{sec:code})؛ وناظمة إيقاع في اليقظة ومتلقيًا صامتًا في النوم. أيّ جهاز سينتج تقرره القنوات الأيونية في الغشاء وقنوات البث التي تضبطها.

\section{أربعة أجهزة لم ترد في الكتاب بعد}

\subsection{الرنّان}
\label{sec:resonator}

لبعض العصبونات تيار بطيء يعاكس أي تغير في الجهد: إذا ارتفع الجهد شدّه التيار إلى الأسفل، وإذا انخفض شدّه إلى الأعلى، لكن بتأخر. هذا التيار يتصرف عند مهندس الإلكترونيات كالمحاثة، ويشكّل مع سعة الغشاء دارة رنين. يستجيب العصبون أقوى ما يستجيب لمدخل بتردد محدد، عادةً من بضعة هرتزات إلى عشرة أو عشرين، ويستجيب أضعف للأبطأ وللأسرع~\cite{HutcheonYarom2000}. إنه مرشح تمرير نطاقي في خلية واحدة: يختار من المدخل المضجوج إيقاعًا بتردده، كالإيقاع المحلي في الفصل~\ref{sec:gaba} مثلًا.

\subsection{المكامل غير التسريبي}
\label{sec:perfectint}

المكامل التسريبي لا يتذكر المدخل إلا زمنًا قدره $\tau$، أي عشرات الأجزاء من الألف من الثانية. لكن العين، كي تثبت، تحتاج إلى أن تتذكر موضعها ثوانيَ: أمر ”الدوران“ يأتي سرعةً قصيرة الأمد، والعين يجب أن تُمسَك في الموضع الجديد. تحل الشبكة ذلك بالتغذية الراجعة نفسها التي بُنيت بها الذاكرة في الفصل~\ref{sec:glutcomputer}. إذا كانت مجموعة من العصبونات ثابتها الزمني $\tau$ تثير نفسها بوزن $w$، فإن $\tau\,\dd r/\dd t=-r+w\,r+I$، والثابت الزمني للمجموعة يساوي
\begin{empheq}[box=\keybox]{equation}
\tau_{\text{eff}} \;=\; \frac{\tau}{1-w} .
\label{eq:perfectint}
\end{empheq}
عند $\tau=0.1\,\text{s}$ و$w=0.995$ نحصل على $20\,\text{s}$: مكامل شبه مثالي من عصبونات ينسى كل منها وحده في عُشر ثانية~\cite{Seung1996}. والثمن هو الضبط الدقيق: إذا زاد $w$ قليلًا على الواحد انجرف المكامل إلى الإشباع، وإذا نقص قليلًا نسي بسرعة. لذلك يحتاج هذا المكامل إلى ضبط دائم بالتعلم، كالذي وُصف في الفصل~\ref{sec:motorlearning}.

\subsection{العصبون ثنائي الاستقرار}
\label{sec:bistable}

في الفصلين~\ref{sec:glutcomputer} و\ref{sec:gaba} جُمع القادح من مجموعة عصبونات. لكن القادح قد يكون في خلية واحدة أيضًا. لبعض العصبونات تيار إذا اشتغل مرة حافظ على الإثارة بنفسه: مدخل قصير ينقل العصبون إلى عمل طويل، هو \emph{الهضبة}، والتثبيط يعيده إلى الراحة. هذا قادح شميت: للعصبون حالتان مستقرتان وبينهما تخلّف. هكذا تتصرف مثلًا عصبونات \nt{ACET} التي تتحكم في العضلات، وتشغّل هذه الخاصيةَ قناةُ بث: تظهر الهضبة حين يصل السيروتونين إلى العصبون~\cite{Hounsgaard1988}. الخلية نفسها ماسك مع \nt{SERO}، ومكامل عادي من دونه.

\subsection{المكاملات والرنّانات}

تقسم النظرية الخلايا القابلة للإثارة إلى صنفين~\cite{Izhikevich2007}. \emph{المكاملات} تشبه مذبذبًا متحكمًا فيه بالجهد يبدأ من الصفر: فوق العتبة بقليل يعمل العصبون ببطء كيفما شئنا، ويرتفع التردد بسلاسة مع المدخل. مثل هذا العصبون يجمع كل ما وصل، وينقل مقدار المدخل بالتردد نقلًا جيدًا. \emph{والرنّانات} تشبه مذبذبًا له تردد أدنى: لا تستطيع العمل ببطء، فعند المدخل الحدّي تُصدر فورًا نبضات بتردد منتهٍ وتفضّل المدخل الذي بترددها. الأولى هي الأجهزة الطبيعية لترميز التردد في الفصل~\ref{sec:code}، والثانية لترميز التوقيت.

\begin{keyblock}
\begin{proposition}[عصبون واحد، أجهزة كثيرة]
\label{prop:eetypes}
تعمل العصبونات، بحسب وظيفتها الكهربائية، مكاملاتٍ تسريبية وغير تسريبية، وكواشف تزامن، ومرشحات تمرير عالٍ ونطاقي، ومحوّلات جهد إلى تردد، ومقوّمات، ومقسِّمات، ومولّدات، وخطوط تأخير، وقوادح. وأيّ جهاز ينتج من خلية بعينها تحدده قنواتها الأيونية وقنوات البث التي تضبطها. الأوضاع نفسها مُقاسة؛ أما مدى استعمال الدماغ لهذه إعادة التهيئة في الحساب ففرضية في الغالب.
\end{proposition}
\end{keyblock}

يشبه هذا عند المهندس المصفوفة المنطقية القابلة للبرمجة: العتاد واحد، والجهاز تحدده التهيئة. غير أن التهيئة هنا لا تتغير عند البرمجة، بل أثناء العمل، بقنوات البث البطيئة التي تحدثنا عنها في الفصول~\babelsublr{\ref{sec:serocomputer}--\ref{sec:acet}}.

\section{تمارين}

\begin{exercise}[\lvlA]
\label{ex:18:1}
\emph{تفاوت المكامل غير التسريبي.} عصبونات بـ$\tau=0.1$ ثانية تحفظ موضع العين وفق~\eqref{eq:perfectint} مدة $T=1$ ثانية بخطأ لا يتجاوز $5\%$ في الاتجاهين. (أ)~أوجد المدى المسموح لـ$w$. (ب)~$w$ مجموع أوزان $K$ مشبكًا، لكل منها خطأ مستقل قدره $10\%$. عند أي $K$ يقع تشتت المجموع ضمن التفاوت؟
\end{exercise}

\begin{exercise}[\lvlB]
\label{ex:18:2}
\emph{الرنّان دارةً.} لنصف التيار البطيء في القسم~\ref{sec:resonator} بالمتغير $W$: $C\,\dd V/\dd t=-g_LV-g_wW+I$، $\tau_w\,\dd W/\dd t=V-W$. بيّن أن هذه دارة $RC$ مع فرع موازٍ $R_w=1/g_w$، $L_w=\tau_w/g_w$، وأوجد عند $C/g_L=10\ms$ و$\tau_w=100\ms$ و$\alpha=g_w/g_L=4$ تردد الرنين والمكسب $|Z|$ عنده على $|Z(0)|$.
\end{exercise}

\begin{exercise}[\lvlB]
\label{ex:18:3}
\emph{كيف يبدأ المكامل العمل.} (أ)~عصبون $\tau\,\dd V/\dd t=-V+\mu$، $\tau=10\ms$، بعتبة $\theta$ وإعادة ضبط إلى الصفر. عند أي $\mu/\theta-1$ يساوي التردد $1\,\text{Hz}$؟ (ب)~ما التردد الذي يعطيه النموذج $\tau\,\dd v/\dd t=v^2+\mu$ (النبضة ذهاب $v$ إلى $+\infty$، وإعادة الضبط إلى $-\infty$) عند $\mu=0.01$؟
\end{exercise}

\begin{exercise}[\lvlB]
\label{ex:18:4}
\emph{نمذجة: المكامل والرنّان.} نموذج التمرين~\ref{ex:18:2} مع $g_L=1$، $\tau_m=10\ms$، $\tau_w=100\ms$، عتبة $1$ وإعادة ضبط $V\to0$ ($W$ لا يتغير) يتلقى $\mu=1.5\sin2\pi ft$، $f=1\text{--}40\,\text{Hz}$. في أي نطاق ترددي يُصدر العصبون نبضات عند $\alpha=0$ وعند $\alpha=4$؟ قارن بالشرط $1.5\,|Z(f)|>1$.
\end{exercise}

\begin{exercise}[\lvlB]
\label{ex:18:5}
\emph{ماسك من خلية واحدة.} عصبون بتيار هضبة: $\tau\,\dd r/\dd t=-r+F(wr+I)$، $F(x)=\min(1,\max(0,x))$، $\tau=10\ms$، $w=1.5$، والخلفية $I=-0.2$. (أ)~ما أقصر نبضة $I=0.3$ تنقل العصبون من $r=0$ إلى الهضبة؟ (ب)~ما أصغر سعة $B$ لنبضة طويلة $I=-0.2-B$ تعيده إلى الراحة؟
\end{exercise}

\begin{exercise}[\lvlB]
\label{ex:18:6}
\emph{الضبط الذاتي للمكامل.} عند تثبيت النظر يكون المدخل $I=0$، ويعطي مستشعر الانزلاق سرعة الانجراف $e=\dd r/\dd t$، و$\dd w/\dd t=-\eta\,e\,r$. عند $w=0.95$، $\tau=0.1$ ثانية، $\langle r^2\rangle=1$ أوجد $\eta$ الذي يُدخل $|1-w|$ في تفاوت التمرين~\ref{ex:18:1} خلال $60$ ثانية من التثبيت. ماذا يفسد لو جرى التعلم أثناء دوران العين أيضًا؟
\end{exercise}

\begin{exercise}[\lvlC]
\label{ex:18:7}
\emph{لماذا تُعاد تهيئة الجهاز؟} قناة بث تبدّل $\alpha$ في نموذج التمرين~\ref{ex:18:2} بين $0$ و$4$. كم ضعفًا يتفوق الرنّان على المكامل في نسبة الإشارة إلى الضجيج لجيب ضعيف في ضجيج تيار أبيض عند $11\,\text{Hz}$ وعند $1\,\text{Hz}$؟ اقترح مسألة يكون فيها تبديل الوضع نفسه هو المفيد.
\end{exercise}
